\section{Does this carry over to multiple objectives?}
\label{sec:mobo}

Everything above is single-objective, while the benchmark table the whole
question comes from is measured in hypervolume. So the last experiment puts
the same contrast into a two-objective problem and uses the paper's own
machinery --- independent GPs, \texttt{qLogEHVI}, the same solvers --- rather
than anything built for this note.

\paragraph{The structure.} One objective carries the map under test; the other
is held fixed.
\[
x \longrightarrow
\begin{cases}
h_0 \longrightarrow g_1 \longrightarrow f_1 & \text{$g_1$ varies} \\
h_1,h_2,h_3,h_4 \longrightarrow g_2 \longrightarrow f_2 & \text{$g_2$ fixed (Ackley)}
\end{cases}
\]
The composite model observes all five intermediates; the direct model sees only
$(f_1,f_2)$. The input dimension, the five inner functions, $g_2$ and the
initial designs are identical across conditions, so any difference between rows
is attributable to the structure of $g_1$ alone.

Two details matter for the comparison to mean anything. Each objective is
normalised on a fixed reference set, because hypervolume is not scale-free and
$z^2$ and $z^3$ otherwise span very different ranges. And all conditions share
\emph{one} reference point --- taken as the element-wise worst across them ---
since letting each pick its own would make the hypervolumes incommensurable,
which is exactly the comparison being made.

\paragraph{Conditions.} The same three regimes as the single-objective study:

\begin{center}
\begin{tabular}{lll}
\toprule
$g_1$ & & what changes \\
\midrule
$z$   & linear & the direct model is \emph{correct}; nothing is lost \\
$z^3$, $\sinh z$ & bijective, nonlinear & wrong model shape, nothing lost \\
$z^2$ & not bijective & wrong model shape \emph{and} the sign of $h_0$ is lost \\
\bottomrule
\end{tabular}
\end{center}

\noindent So linear~$\to$~bijective isolates misspecification, and
bijective~$\to$~square adds information loss at comparable nonlinearity.

\paragraph{Result.} Four function draws $\times$ eight seeds, $40$ adaptive
evaluations, $32$ paired runs per condition.

\begin{center}
\begin{tabular}{llrrrr}
\toprule
$g_1$ & & HV direct & HV composite & Gain (\%) & median \\
\midrule
$z$    & linear         & $4.338$ & $5.046$ & $+17.3\pm2.1$ & $+15.6$ \\
$z^3$  & bijective      & $5.033$ & $6.267$ & $+30.8\pm6.4$ & $+22.6$ \\
$z^2$  & not bijective  & $3.935$ & $4.905$ & $+25.8\pm2.2$ & $+24.8$ \\
\bottomrule
\end{tabular}
\end{center}

\noindent Composite wins in every condition and on $26$ to $31$ of $32$ runs,
which is the first thing worth saying: the advantage is large and robust in
hypervolume, as in the benchmark table.

\paragraph{But the channels do not separate here.} Pairing by draw and seed:

\begin{center}
\begin{tabular}{lrrr}
\toprule
Contrast & Mean & Median & Sign test \\
\midrule
$z \to z^3$ \quad (misspecification only) & $+13.4\pm6.6$ & $-0.7$ & $16/32$, $p=1.00$ \\
$z^3 \to z^2$ \quad (adds information loss) & $-5.0\pm7.3$ & $+6.5$ & $19/32$, $p=0.38$ \\
$z \to z^2$ \quad (both channels) & $+8.5\pm3.2$ & $+7.4$ & $21/32$, $p=0.11$ \\
\bottomrule
\end{tabular}
\end{center}

\noindent The mean and the median disagree on the first row, and the sign test
reads an exact coin flip. The cause is $z^3$: its gain runs from $-29\%$ to
$+131\%$, so a few seeds where the direct GP fits badly carry the mean well
above the median ($30.8$ against $22.6$). The ordering of the medians
($15.6 \to 22.6 \to 24.8$) is the one predicted, but no individual contrast is
significant.

\paragraph{Why hypervolume is a blunter instrument.} The single-objective study
resolved these channels because CRPS is a per-point score averaged over $512$
held-out points. Hypervolume is one number per run, produced by a stochastic
forty-step optimisation, with a heavy right tail. The standard errors here are
roughly five times worse relative to the effect. Resolving the
bijective-to-square step at this noise level needs about $13\times$ more runs,
roughly $400$ per condition.

That is a statement about the \emph{measurement}, not about the effect. The
exact decomposition of Section~\ref{sec:channels} does not depend on any of
this: it is computed, not estimated, and it already shows the information
channel to be the smaller of the two. The multi-objective experiment is
consistent with that and simply lacks the power to confirm it independently.

\paragraph{What this adds anyway.} Two things that do not need the contrast to
separate. First, the large composite advantage survives in hypervolume on a
problem where the inner functions are drawn from the surrogate's own prior, so
it is not an artefact of any particular benchmark. Second, the advantage is
substantial \emph{even when $g_1$ is linear} ($+17.3\%$, $31/32$ runs) --- a
case where the proposition guarantees that observing $h$ conveys no information
at all. Whatever is producing that gain, it is not information.
